\documentclass[12pt,a4paper]{article}

% Paquetes esenciales
\usepackage[utf8]{inputenc}
\usepackage[spanish]{babel}
\usepackage[T1]{fontenc}
\usepackage{geometry}
\usepackage{xcolor}
\usepackage{booktabs}
\usepackage{enumitem}

% Configuración de márgenes
\geometry{left=2.5cm, right=2.5cm, top=2.5cm, bottom=2.5cm}

% Configuración de colores para títulos
\definecolor{darkblue}{RGB}{0, 51, 102}
\usepackage{titlesec}
\titleformat{\section}{\color{darkblue}\normalfont\Large\bfseries}{\thesection}{1em}{}
\titleformat{\subsection}{\color{darkblue}\normalfont\large\bfseries}{\thesubsection}{1em}{}

\title{\textbf{Análisis Estratégico: Rendimiento de Ventas - Dataset Superstore}}
\author{Luis Fernando Murillo Corral}
\date{\today}

\begin{document}

\maketitle

\section{Objetivos del Análisis}
El presente reporte tiene como finalidad evaluar el rendimiento comercial de la empresa, identificando áreas de oportunidad para optimizar la rentabilidad. Los objetivos principales son:
\begin{itemize}[noitemsep]
    \item Identificar los sectores con mayor volumen de ventas.
    \item Analizar las categorías con mayor generación de ganancias durante el periodo 2016-2018.
    \item Proponer estrategias de mejora en sectores con bajo desempeño.
\end{itemize}

\section{Análisis de Situación}

\subsection{Sectores con mayor volumen de ventas}
Tras el análisis de datos, se identifica que la categoría de \textbf{Tecnología} es el motor principal de ventas. En particular, la modalidad de envío \textit{Standard Class} concentra el mayor volumen de transacciones, lo que sugiere una alta demanda de productos tecnológicos en envíos regulares, representando un porcentaje significativo del volumen total.

\subsection{Análisis de rentabilidad (2016-2018)}
El año 2018 marcó el pico de rentabilidad en el periodo analizado. Se observa una contracción temporal en 2016 dentro de las categorías de Tecnología y Office Supplies. Un hallazgo clave es la estacionalidad: los meses de noviembre y diciembre presentan los mayores niveles de ingresos. Este fenómeno responde tanto a la temporada de festividades como a los posibles cierres presupuestarios corporativos de fin de año.

\subsection{Sectores con oportunidad de mejora}
La categoría de \textit{Furniture} (Mobiliario) mantiene un crecimiento constante pero lineal. En contraste, la categoría de \textit{Tecnología} posee un campo de innovación más amplio, lo que la posiciona como la candidata ideal para mayores inversiones de marketing y desarrollo estratégico.

\section{Propuestas de Mejora y Estrategias}

\begin{itemize}
    \item \textbf{Optimización de campañas estacionales:} Dado que noviembre y diciembre son meses críticos, es necesario revisar la caída de ventas observada entre ambos meses en 2018. Se propone una estrategia de marketing segmentada que mantenga el \textit{momentum} de noviembre hacia la última semana de diciembre.
    
    \item \textbf{Expansión estratégica en California:} Siendo California el estado con mayor generación de ingresos, se recomienda implementar un plan de conexión logística con zonas aledañas para fomentar la diversificación de rutas y maximizar el alcance de clientes actuales.
    
    \item \textbf{Estrategia de precios "Señuelo":} Para incentivar la compra de productos de \textit{First Class}, se propone ajustar los precios de \textit{Second Class} manteniendo una diferencia menor al 10\% respecto a la primera. Esta técnica de arquitectura de precios busca que el cliente perciba un mayor valor en la clase superior. 
    \textit{Nota: Esta medida será evaluada mediante un test A/B en regiones seleccionadas para medir la elasticidad de la demanda.}
\end{itemize}

\end{document}
